\section{PROBLEM FORMULATION}
\label{sec-problem-formulation}
This section formalizes resilient restoration under adversarial recovery knowledge. The formulation separates four issues that were conflated in the original manuscript: partial observability, evidence-dependent action admissibility, worst-case recovery performance, and the security objective of bounding the propagation of erroneous knowledge. The attacker is modeled through the observation and evidence processes defined in Section~\ref{sec-system-model}, rather than as a cooperative decision agent.

\subsection{Constrained Dec-POMDP under Adversarial Knowledge}
\label{subsec-decpomdp}
Let $\mathcal{I}=\{1,\ldots,N\}$ denote the set of restoration agents. Because each agent observes only a local part of the PCPS and the recovery knowledge state summarizes time-dependent evidence history, we model restoration as a constrained decentralized partially observable Markov decision process (C-Dec-POMDP)
\begin{align}
\label{eq-cdecpomdp}
\mathcal{M}_{\rm rec}=\left\langle
\mathcal{I},\mathcal{X},\{\mathbb{U}_i\}_{i\in\mathcal{I}},
P,\{\mathcal{O}_i\}_{i\in\mathcal{I}},O,R,\mathcal{C},\gamma
\right\rangle,
\end{align}
where $\mathcal{X}$ is the latent cross-layer state space, $\mathbb{U}_i$ is the action space of agent $i$, $P(x_{t+1}\mid x_t,\mathbf{a}_t)$ is the state-transition kernel, $\mathcal{O}_i$ is the local observation space, $O$ is the observation model, $R$ is the cooperative restoration reward, $\mathcal{C}$ is the constraint set, and $\gamma\in(0,1)$ is the discount factor. We use $\mathbb{U}_i$ for action spaces to avoid overloading $\mathcal{A}_t$, which denotes dynamic KG assertions in Section~\ref{subsec-recovery-knowledge-plane}.

At time $t$, agent $i$ receives a local observation $o_{i,t}$ and has access to the task-relevant portion of the recovery belief state $\mathcal{B}_t$. Its decision information is written abstractly as
\begin{align}
\label{eq-information-state}
\iota_{i,t}=\big(\tau_{i,t},\mathcal{B}^{i}_t\big),
\end{align}
where $\tau_{i,t}=(o_{i,0},a_{i,0},\ldots,o_{i,t})$ is the local action--observation history and $\mathcal{B}^{i}_t\subseteq\mathcal{B}_t$ contains the recovery beliefs relevant to the agent's current task. The concrete extraction and encoding of $\mathcal{B}^{i}_t$ are implementation choices deferred to Section~V. A decentralized policy selects
\begin{align}
\label{eq-decentralized-policy}
a_{i,t}\sim\pi_i(\cdot\mid\iota_{i,t}),
\end{align}
and the joint candidate action is $\mathbf{a}_t=(a_{1,t},\ldots,a_{N,t})$.

The adversary affects decision making indirectly by manipulating admissible evidence streams within the budget $\mathcal{B}_{A}$ in Eq.~\ref{eq-attack-budget}. Thus the observation and knowledge processes are attack dependent, whereas the physical transition kernel remains governed by the PCPS dynamics in Eq.~\ref{eq-system-transition}. This distinction is important: the security problem is not merely to tolerate exogenous state disturbances, but to prevent corrupted observations from becoming trusted recovery knowledge and then executable control authority.

\subsection{Risk-Aware Action Admissibility}
\label{subsec-risk-admissibility}
We partition the restoration constraints into invariant engineering constraints and evidence-dependent prerequisites,
\begin{align}
\label{eq-constraint-partition}
\mathcal{C}=\mathcal{C}^{\rm inv}\cup\mathcal{C}^{\rm evid}_t.
\end{align}
The set $\mathcal{C}^{\rm inv}$ contains requirements whose validity does not depend on online attack diagnosis, such as voltage and thermal limits, device ratings, power-balance relations, protection interlocks, and radiality. By contrast, $\mathcal{C}^{\rm evid}_t$ contains runtime conditions inferred from evidence, such as whether a controller is available, whether a command route is reachable, or whether a measurement chain is sufficiently trustworthy for a specific operation.

For a candidate joint action $\mathbf{a}$, let
\begin{align}
\mathcal{Q}(\mathbf{a})=\{q_1,\ldots,q_m\}
\end{align}
be the dynamic assertions that must hold for its cyber-physical execution. Instead of converting each KG assertion into a binary rule, we assign a risk budget $\delta_q\in(0,1)$ and require
\begin{align}
\label{eq-belief-prerequisite}
b_t(q)\ge 1-\delta_q,
\qquad \forall q\in\mathcal{Q}(\mathbf{a}).
\end{align}
The corresponding belief-admissible set is
\begin{align}
\label{eq-belief-admissible-set}
\mathbb{U}^{B}_t=
\left\{\mathbf{a}\ \middle|\
 b_t(q)\ge1-\delta_q,
\ \forall q\in\mathcal{Q}(\mathbf{a})
\right\}.
\end{align}
Equation~\ref{eq-belief-admissible-set} makes uncertainty explicit: insufficient evidence cannot be interpreted as permission. The aggregate semantic risk budget for action $\mathbf{a}$ is constrained by
\begin{align}
\label{eq-semantic-risk-budget}
\sum_{q\in\mathcal{Q}(\mathbf{a})}\delta_q\le\bar{\delta}_{B},
\end{align}
where $\bar{\delta}_{B}$ is selected according to the criticality of the restoration operation.

Invariant safety is checked independently of the belief gate. Let $\widehat{\mathcal{X}}_t$ be a conservative set of physical states consistent with the currently trusted measurements and uncertainty bounds. We define
\begin{align}
\label{eq-invariant-admissible-set}
\mathbb{U}^{\rm inv}_t=
\Big\{\mathbf{a}\ \Big|\
&g_k(x,\mathbf{a})\le0,\ \forall k,\\
&h_l(x,\mathbf{a})=0,\ \forall l,
\ \forall x\in\widehat{\mathcal{X}}_t
\Big\},\nonumber
\end{align}
where $g_k$ and $h_l$ represent inequality and equality/structural constraints, respectively. The final admissible action set is
\begin{align}
\label{eq-final-admissible-set}
\mathbb{U}^{\rm adm}_t
=\mathbb{U}^{B}_t\cap\mathbb{U}^{\rm inv}_t.
\end{align}
A learned surrogate may later reduce the number of candidates sent to exact verification, but it cannot enlarge $\mathbb{U}^{\rm inv}_t$ and therefore has no independent authorization authority. If no nontrivial action is admissible, the controller executes a certified fallback action $\mathbf{a}^{\rm fb}_t\in\mathbb{U}^{\rm inv}_t$, such as maintaining isolation, retaining local control, or shedding a bounded amount of load.

\subsection{Robust Recovery Objective}
\label{subsec-robust-objective}
Let $\mathcal{L}$ be the set of loads, $P_l$ the nominal demand of load $l$, $w_l>0$ its service priority, and $y_{l,t}\in[0,1]$ its restored fraction at time $t$. The weighted restored-service ratio is
\begin{align}
\label{eq-weighted-recovery-ratio}
\eta_t=
\frac{\sum_{l\in\mathcal{L}}w_lP_l y_{l,t}}
{\sum_{l\in\mathcal{L}}w_lP_l}.
\end{align}
We use safety as a hard constraint rather than allowing a large restoration reward to compensate for an unsafe operation. The stage reward therefore captures service recovery and operational costs,
\begin{align}
\label{eq-robust-reward}
R_t=
\lambda_{\eta}\eta_t
-\lambda_{s}C^{\rm shed}_t
-\lambda_{o}C^{\rm op}_t
-\lambda_{c}C^{\rm cyber}_t
-\lambda_{f}C^{\rm fb}_t,
\end{align}
where $C^{\rm shed}_t$, $C^{\rm op}_t$, $C^{\rm cyber}_t$, and $C^{\rm fb}_t$ denote critical-load curtailment, switching/dispatch cost, cyber-resource recovery cost, and fallback-service loss, respectively.

Let $\mathfrak{A}(\mathcal{B}_{A})$ denote the set of adversaries satisfying the capability and budget assumptions in Section~\ref{subsec-threat-model}. We seek a recovery policy and online belief-update rule that maximize worst-case expected recovery performance,
\begin{align}
\label{eq-minimax-objective}
\max_{\Pi,F_B}\ \inf_{\mathscr{A}\in\mathfrak{A}(\mathcal{B}_{A})}
\mathbb{E}_{\mathscr{A}}
\left[\sum_{t=0}^{T}\gamma^t R_t\right],
\end{align}
where $\Pi=\{\pi_i\}_{i\in\mathcal{I}}$ and
$F_B$ denotes the online update mapping that produces $\mathcal{B}_{t+1}$ from the previous belief state and newly received evidence. The optimization in Eq.~\ref{eq-minimax-objective} is subject to
\begin{align}
\label{eq-execution-constraint}
\mathbf{a}^{\rm exec}_t\in\mathbb{U}^{\rm adm}_t,
\end{align}
and to the physical safety and recovery requirements
\begin{align}
\label{eq-safety-probability}
\Pr\left(x_t\in\mathcal{X}_{\rm unsafe}\right)&\le\delta_x,
\qquad 0\le t\le T,\\
\label{eq-recovery-liveness}
\Pr\left(\exists t\le T_{\max}:x_t\in\mathcal{X}_{\rm tar}\right)&\ge1-\delta_r.
\end{align}
Here $\mathcal{X}_{\rm unsafe}$ and $\mathcal{X}_{\rm tar}$ denote the unsafe and target recovery sets, respectively. Equation~\ref{eq-recovery-liveness} prevents an overly conservative knowledge gate from achieving safety simply by refusing all useful restoration actions.

\subsection{Security Objectives: Bounding Knowledge Propagation}
\label{subsec-security-objectives}
The distinguishing security objective of this work is to bound how far adversarial evidence can propagate through the recovery loop. We formalize this requirement at the knowledge, action, correction, and policy levels. These properties are design targets in this section; Section~V provides the mechanisms and the later theoretical analysis establishes sufficient conditions for them.

\emph{1) Bounded belief distortion.}
Let $\mathcal{B}^{0}_t$ denote the counterfactual recovery belief state generated from uncompromised evidence under the same physical trajectory, and let $\mathcal{B}^{A}_t$ denote the belief state under an admissible attack. For the task-relevant assertion set $\mathcal{Q}_t$, define
\begin{align}
\label{eq-belief-distance}
D_B(\mathcal{B}^{A}_t,\mathcal{B}^{0}_t)
=\sum_{q\in\mathcal{Q}_t}\omega_q
\left|b^{A}_t(q)-b^{0}_t(q)\right|,
\qquad \sum_q\omega_q=1.
\end{align}
The first objective is the existence of a finite attack-budget-dependent bound
\begin{align}
\label{eq-bounded-belief-distortion}
\sup_{t\le T}
D_B(\mathcal{B}^{A}_t,\mathcal{B}^{0}_t)
\le \varepsilon_B(\mathcal{B}_{A}).
\end{align}
This rules out unbounded accumulation of influence from a bounded number of compromised provenance groups.

\emph{2) Bounded false authorization.}
Let $\mathbb{U}^{\star}_t$ denote the action set that would be semantically admissible if the ground-truth values $q_t^{\star}$ of all prerequisites were known. A false authorization occurs when an action is admitted by the online belief state but is not contained in $\mathbb{U}^{\star}_t$. We require
\begin{align}
\label{eq-false-authorization-bound}
\Pr\left(
\mathbf{a}^{\rm exec}_t\in
\mathbb{U}^{\rm adm}_t\setminus\mathbb{U}^{\star}_t
\right)
\le \delta_{\rm auth}.
\end{align}
This objective formalizes the desired barrier between erroneous knowledge and field control authority.

\emph{3) Finite correction after poisoning.}
Suppose an attack affecting assertion $q$ ceases at time $t_a$ and sufficient uncompromised evidence resumes. For tolerance $\epsilon_{\rm corr}>0$, define the correction time
\begin{align}
\label{eq-correction-time}
T_{\rm corr}(q)=
\inf\left\{k\ge0:\left|b^{A}_{t_a+k}(q)-b^{0}_{t_a+k}(q)\right|\le\epsilon_{\rm corr}\right\}.
\end{align}
We seek conditions under which
\begin{align}
\label{eq-finite-correction}
T_{\rm corr}(q)\le\overline{T}_{\rm corr}
\end{align}
for all task-critical assertions whose honest evidence becomes observable again. This prevents stale or poisoned recovery facts from persisting indefinitely after the attack ends.

\emph{4) Bounded knowledge-to-policy propagation.}
Let $\Pi^{A}_t$ and $\Pi^{0}_t$ denote online policies adapted from attacked and uncompromised recovery knowledge, respectively. We measure adversarially induced policy deviation by
\begin{align}
\label{eq-policy-distance}
D_{\pi}(\Pi^{A}_t,\Pi^{0}_t)
=\frac{1}{N}\sum_{i=1}^{N}
\mathbb{E}_{\iota\sim\mu_t}
\left[
D_{\rm KL}
\left(
\pi^{A}_i(\cdot\mid\iota)
\,\|\,
\pi^{0}_i(\cdot\mid\iota)
\right)
\right],
\end{align}
where $\mu_t$ is the task-relevant information-state distribution. The required propagation bound is
\begin{align}
\label{eq-policy-propagation-bound}
D_{\pi}(\Pi^{A}_t,\Pi^{0}_t)
\le\varepsilon_{\pi}(\mathcal{B}_{A}).
\end{align}
The mechanism in Section~V enforces this objective by separating fast knowledge updates from slower trust-gated policy adaptation, so that a newly inserted or uncertain KG assertion cannot immediately alter the deployed recovery policy.

Together, Eqs.~\ref{eq-bounded-belief-distortion}--\ref{eq-policy-propagation-bound} define the core security problem addressed by this paper: under incomplete, stale, conflicting, and partially adversarial recovery evidence, maintain an online knowledge state whose error remains bounded, prevent that error from directly authorizing unsafe recovery actions, correct poisoned knowledge within finite time once trustworthy evidence returns, and limit the resulting deviation of online recovery policies.